This report established a rigorous, reproducible performance baseline for RF
fingerprinting on a 12-device dataset. A complete DL pipeline was evaluated under a rigorous, group-aware protocol;
all models achieve ADR in the 0.57--0.68 range,
confirming that RF transient emissions carry sufficient device-discriminating
information to support PLA above chance, while highlighting per-classifier
sample count as the binding constraint. Addressing this --- through larger
datasets or improved augmentation strategies --- is the most direct path
toward operationally useful PLA. Future directions are listed in
Section~\ref{sec:contributions}.

Two results have direct implications for practical deployment. First, the
high seed-to-seed variance of the BiGRU ($\Delta$ADR up to 0.37 across seeds
for a single trial) means the system is not yet reliable enough for deployment
at fixed thresholds; an ensemble or calibration step is needed before
per-device classifiers can operate at a stable operating point. Second, the
multi-class probe --- which achieves substantially higher relative gain over
chance despite identical data --- suggests that joint training across all
devices is the right inductive bias for a dataset of this size, and warrants
prioritisation over further tuning of per-device binary classifiers.
